% !TeX root = ../../Book1.tex
% 纲要标题：### 13.1 分式域
% 纲要位置：计划纲要.md，第 472 行。

\section{分式域}
\label{b1:sec:1301}

\cref{b1:prop:fraction-field-basic}已经完整构造整环的分式域，并验证等价关系、运算和域公理。本节沿用这个对象，研究它与其他域的关系：在已有的域中如何识别分式域，从整环出发的同态又何时能够延拓到分式域？这些问题将引出泛性质，并为下一节只让指定元素取逆的一般局部化作准备。

% 纲要内容（仅作写作计划，尚非教材正文）：
%
% * 分式域的回顾与例子；完整构造已在第12.2节建立
% * 分式相等与分式域的识别
% * 分式域的泛性质，为一般局部化准备同态延拓的观点
%

\subsection{分式域的回顾与例子}
\label{b1:subsec:1301:01}

设 \(R\) 为交换整环。回顾\cref{b1:prop:fraction-field-basic}，域 \(\operatorname{Frac}(R)\) 的元素由 \(a/b\) 表示，其中 \(a,b\in R\)、\(b\neq0\)，相等的判据是
\begin{equation}\label{b1:eq:fraction-equivalence}
\frac ab=\frac cd\quad\Longleftrightarrow\quad ad=bc.
\end{equation}
这里两个分母都非零，整环的消去律使交叉乘积相等成为完整的判据。在有零因子的一般环中，\(u(ad-bc)=0\) 未必允许消去 \(u\)；但若让 \(u\) 变成单位，交叉乘积的差便会变成零。下一节分式等价关系中的额外因子，正是为了记录这种因取逆而消失的差异。

\ProseParagraphBreak
后面使用的运算公式为
\begin{equation}\label{b1:eq:fraction-operations}
\frac ab+\frac cd=\frac{ad+bc}{bd},
\qquad
\frac ab\cdot\frac cd=\frac{ac}{bd}.
\end{equation}
分式域也可以不指定等价类模型，而用它包含原环的方式来刻画。

\begin{definition}\label{b1:def:fraction-field}
设 \(R\) 为交换整环，\(K\) 为域，\(i:R\rightarrow K\) 为保持单位的单射环同态。如果每个 \(x\in K\) 都能写成
\[
x=i(a)i(b)^{-1}\qquad(a,b\in R,\ b\neq0),
\]
就称域 \(K\) 连同嵌入 \(i\) 为 \(R\) 的分式域。
\end{definition}
已构造的 \(\operatorname{Frac}(R)\) 与嵌入 \(i(a)=a/1\) 满足这个定义。通常通过 \(i\) 把 \(R\) 看成分式域的子环，并将 \(i(a)\) 简记为 \(a\)。本节末将证明：不同模型之间有唯一保持这个嵌入的同构。

\ProseParagraphBreak
例如 \(\operatorname{Frac}(\mathbb Z)=\mathbb Q\)。若 \(k\) 为域，则 \(k[x]\) 的分式域记为 \(k(x)\)，称为 \(k\) 上的\term[youli-hanshu-yu]{有理函数域}[rational function field]；元素是多项式的分式 \(f(x)/g(x)\)，其中 \(g\neq0\)。
\symidx{\(k(x)\)：域 \(k\) 上的一元有理函数域}
这里说的是形式分式，而不是在 \(k\) 上处处有定义的函数。域元素 \(h\in k(x)\) 可以有多个分式表示 \(h=f/g\)；对于指定的一个表示，只有 \(g(a)\ne0\) 时才能在 \(a\in k\) 处直接代入。例如 \(1/x\) 是 \(k(x)\) 的元素，但所写表示不能在 \(x=0\) 处直接代入。若 \(R\) 本来就是域，则 \(a/b\mapsto ab^{-1}\) 给出 \(\operatorname{Frac}(R)\cong R\)，构造并没有加入新的元素。

\subsection{分式相等与分式域的识别}
\label{b1:subsec:1301:02}

分式计算的依据是\eqref{b1:eq:fraction-equivalence}，不必先选定所谓最简分式。一般整环中的最大公因子未必存在，因而不能把“约分到最简”作为分式域定义的一部分。即使分子、分母只有单位作为公因子，也未必能保证两种表示只差同时乘一个单位；\cref{b1:ex:13-13}将在 \(\mathbb Z[\sqrt{-5}]\) 中给出这样的既约分式反例。

\ProseParagraphBreak
在有零因子的环中，分母取逆还可能使自然映射不再单射：若 \(sa=0\) 而要求 \(s\) 可逆，那么在新环中 \(a=s^{-1}(sa)=0\)，即使原环中的 \(a\) 非零也一样。因此下一节的局部化既要加入指定元素的逆，也可能使原来的某些元素变成零。

\ProseParagraphBreak
分式相等不表示分子、分母分别相等。例如在 \(k(x)\) 中，
\[
\frac{x^2-1}{x-1}=x+1,
\]
因为 \(x^2-1=(x-1)(x+1)\)，而多项式 \(x-1\) 非零。这是有理函数域内的等式；左端原表示在 \(x=1\) 处不能直接代入，右端则可以。给定域元素 \(h\in k(x)\)，如果它至少有一种表示 \(h=f/g\) 满足 \(g(a)\ne0\)，就说 \(h\) 在 \(a\) 处正则，并用这一表示求值。这个值不依赖所选的可代入表示：若 \(f/g=f_1/g_1\)，且两个分母在 \(a\) 处都非零，则把 \(fg_1=f_1g\) 在 \(a\) 处求值，再消去非零分母，便得到相同的商值。因此，在一点正则是域元素具有某种可代入表示的性质，不要求它的每一种表示都能直接代入；上式中的元素在 \(1\) 处正则。

\ProseParagraphBreak
若 \(R\) 已经放在一个域 \(L\) 中，每个包含 \(R\) 的子域都被迫包含非零 \(b\in R\) 的逆，进而包含所有 \(ab^{-1}\)。所以，要在 \(L\) 中识别分式域，就应收集这些必需的商，并检查它们自身已经组成子域；这时它便是包含 \(R\) 的最小子域。

\begin{proposition}\label{b1:prop:fraction-in-ambient-field}
设 \(R\) 是域 \(L\) 的子环。则
\[
K=\left\{\,ab^{-1}\in L\mid a,b\in R,\ b\neq0\,\right\}
\]
是 \(L\) 中包含 \(R\) 的最小子域，且 \(a/b\mapsto ab^{-1}\) 是从 \(\operatorname{Frac}(R)\) 到 \(K\) 的同构。
\end{proposition}
\begin{proof}
对 \(a,b,c,d\in R\)、\(b,d\neq0\)，域 \(L\) 中有
\[
ab^{-1}+cd^{-1}=(ad+bc)(bd)^{-1},
\qquad
(ab^{-1})(cd^{-1})=(ac)(bd)^{-1}.
\]
这些元素仍在 \(K\) 中，且 \(K\) 含 \(0,1\)，在取负和对非零元素取逆下也封闭，所以是子域。包含 \(R\) 的子域必须包含每个 \(b^{-1}\) 及 \(ab^{-1}\)，故包含 \(K\)。最后，\(ab^{-1}=cd^{-1}\) 在 \(L\) 中等价于 \(ad=bc\)；这个等式也在子环 \(R\) 中成立。因此所给映射良定义且单射，分式运算说明它是环同态，\(K\) 的定义保证满射。
\end{proof}
这个命题允许在已有的域中直接寻找分式域。例如把 \(\mathbb Z[i]\) 放在 \(\mathbb C\) 中，其分式域就是 \(\mathbb Q(i)=\{\,a+bi\mid a,b\in\mathbb Q\,\}\)。后一集合在和、积及取负下封闭；非零元素的逆为
\[
(a+bi)^{-1}=\frac{a-bi}{a^2+b^2},
\]
仍具有有理数的实部与虚部，所以它是包含 \(\mathbb Z[i]\) 的域。每个这样的元素又可通分成一个 Gauss 整数除以非零整数，故确为 \(\mathbb Z[i]\) 的分式域。

\subsection{分式域的泛性质}
\label{b1:subsec:1301:03}

把分式代入映射时，分母必须仍可取逆。这给出分式域的泛性质。目标可以是一般交换含幺环 \(A\)：不要求其中每个非零元素都可逆，只要求原环中将被取逆的非零元素的像可逆。这正是下一节局部化泛性质所保留的条件。

\begin{theorem}\label{b1:thm:fraction-field-universal}
设 \(R\) 为交换整环，\(A\) 为交换含幺环，\(f:R\rightarrow A\) 为保持单位的环同态，记自然嵌入为 \(\iota(a)=a/1\)。存在保持单位的环同态
\(\widetilde f:\operatorname{Frac}(R)\rightarrow A\) 满足 \(\widetilde f(a/1)=f(a)\) 对每个 \(a\in R\) 成立，当且仅当每个 \(b\in R\setminus\{0\}\) 的像 \(f(b)\) 都是 \(A\) 的单位。存在时，此同态唯一；对 \(a,b\in R\)、\(b\ne0\)，有
\[
\widetilde f\left(\frac ab\right)=f(a)f(b)^{-1}.
\]
延拓存在时，下面的图交换：
\[
\begin{tikzcd}[column sep=large,row sep=large]
R\arrow[r,"\iota"]\arrow[dr,"f"']
&\operatorname{Frac}(R)\arrow[d,dashed,"\exists!\,\widetilde f"]\\
&A.
\end{tikzcd}
\]
若 \(A\ne0\) 且上述延拓存在，则 \(f\) 和 \(\widetilde f\) 都单射。特别地，若 \(A=L\) 为域而 \(f\) 单射，则延拓存在，并且是域的单射同态。
\end{theorem}
\begin{proof}
若延拓存在，\(b/1\) 的逆元为 \(1/b\)，故 \(f(b)\widetilde f(1/b)=1\)，说明 \(f(b)\) 必须可逆。反过来，假设每个这样的 \(f(b)\) 可逆。若 \(a/b=c/d\)，则 \(ad=bc\)，从而
\(f(a)f(d)=f(b)f(c)\)。乘以 \(f(b)^{-1}f(d)^{-1}\)，得到
\(f(a)f(b)^{-1}=f(c)f(d)^{-1}\)，所给公式因而良定义。由\eqref{b1:eq:fraction-operations}，
\[
\begin{aligned}
\widetilde f\left(\frac ab+\frac cd\right)
&=\bigl(f(a)f(d)+f(b)f(c)\bigr)f(b)^{-1}f(d)^{-1}\\
&=f(a)f(b)^{-1}+f(c)f(d)^{-1},\\
\widetilde f\left(\frac ab\cdot\frac cd\right)
&=f(a)f(c)f(b)^{-1}f(d)^{-1}.
\end{aligned}
\]
故它保持加法、乘法及单位。任何延拓都必须把 \(a/b=(a/1)(b/1)^{-1}\) 映为 \(f(a)f(b)^{-1}\)，所以唯一。若 \(A\ne0\) 且延拓条件满足，假设有 \(0\ne a\in\ker f\)，则 \(f(a)=0\) 必须是单位，但非零环中的零元不可逆，矛盾。因此 \(f\) 单射。再由 \(\widetilde f(a/b)=0\) 得 \(f(a)=0\)，进而 \(a=0\)，故延拓也单射。若 \(A=L\) 为域且 \(f\) 单射，每个非零 \(b\) 的像都非零，因而可逆，延拓条件自动满足。
\end{proof}

例如单射 \(\mathbb Z\hookrightarrow\mathbb R\) 唯一地延拓为 \(\mathbb Q\hookrightarrow\mathbb R\)。目标不必是域：对角映射 \(\mathbb Z\rightarrow\mathbb Q\times\mathbb Q\)、\(a\mapsto(a,a)\) 也将每个非零整数送到单位，故唯一延拓为 \(a/b\mapsto(a/b,a/b)\)，尽管 \(\mathbb Q\times\mathbb Q\) 有零因子。模素数 \(p\) 的商映射 \(\mathbb Z\rightarrow\mathbb F_p\) 则不能延拓到 \(\mathbb Q\)：它的目标非零而核含有非零整数 \(p\)，不满足上面自动单射的必要条件；等价地说，\(p\) 的像为零，不能作逆。这个障碍不能靠另选分式代表来消除。

\begin{corollary}\label{b1:cor:fraction-field-unique}
设 \(R\) 为交换整环。若 \((K,i)\)、\((K',i')\) 都是 \(R\) 的分式域，则存在唯一保持指定 \(R\) 嵌入的域同构 \(\varphi:K\rightarrow K'\)，即 \(\varphi\circ i=i'\)。这样的同构记作 \((K,i)\cong_R(K',i')\)。
\end{corollary}
\begin{proof}
由\cref{b1:prop:fraction-in-ambient-field}，映射
\[
u:\operatorname{Frac}(R)\xrightarrow{\sim}K,
\qquad
v:\operatorname{Frac}(R)\xrightarrow{\sim}K',
\]
其中 \(u(a/b)=i(a)i(b)^{-1}\)、\(v(a/b)=i'(a)i'(b)^{-1}\)，都是同构：按分式域的定义，\(K\)、\(K'\) 中每个元素都可写成相应的商。取
\[
\varphi=v\circ u^{-1}:K\xrightarrow{\sim}K',
\]
便有 \(\varphi\circ i=i'\)。若另有这样的同态，则它在每个 \(i(a)i(b)^{-1}\) 上的值必为 \(i'(a)i'(b)^{-1}\)，所以与 \(\varphi\) 一致。
\end{proof}
因此，具体等价类只是构造模型；决定分式域的是它对 \(R\) 的包含关系和分母取逆的要求。强调“保持 \(R\) 的嵌入”也很必要：抽象域之间可能有多个同构，唯一性指满足指定条件的那一个。
